% !TeX program = xelatex
% =====================================================================
%  人工智能与数学 —— 课程作业报告
%  题目：CRP-Agent：面向极值图论不等式的“猜想—反驳—证明”智能体
%
%  编译方式：
%      xelatex main.tex
%      xelatex main.tex        （第二遍用于生成目录与交叉引用）
%  或在 report/ 目录下执行： latexmk -xelatex main.tex
% =====================================================================
\documentclass[12pt, a4paper]{ctexart}

% ------------------------- 页面与版式 -------------------------
\usepackage{geometry}
\geometry{left=2.6cm, right=2.6cm, top=2.7cm, bottom=2.7cm}

% ------------------------- 数学 -------------------------
\usepackage{amsmath, amssymb, amsthm, mathtools, bm}

% ------------------------- 图表与排版 -------------------------
\usepackage{graphicx}
\usepackage{booktabs}
\usepackage{longtable}
\usepackage{tabularx}
\usepackage{float}
\usepackage{enumitem}
\usepackage{xcolor}
\usepackage{caption}
\usepackage{titlesec}
\usepackage{fancyhdr}
\usepackage{microtype}

% ------------------------- 绘图 -------------------------
\usepackage{tikz}
\usetikzlibrary{arrows.meta, positioning, shapes.geometric, fit, backgrounds, calc}

% ------------------------- 强调框 -------------------------
\usepackage{tcolorbox}
\tcbuselibrary{skins, breakable}

% ------------------------- 超链接（放最后） -------------------------
\usepackage[colorlinks=true,
            linkcolor=blue!45!black,
            citecolor=blue!45!black,
            urlcolor=blue!45!black,
            bookmarksnumbered=true]{hyperref}

% =====================================================================
%  自定义
% =====================================================================
\definecolor{accent}{RGB}{20, 80, 140}
\definecolor{accentlight}{RGB}{235, 243, 250}
\definecolor{warnlight}{RGB}{252, 245, 235}

% 定理类环境：共用计数器，按节编号
\theoremstyle{definition}
\newtheorem{definition}{定义}[section]
\newtheorem{problem}{问题}[section]
\newtheorem{example}{例}[section]
\newtheorem{remark}{注}[section]

\theoremstyle{plain}
\newtheorem{theorem}{定理}[section]
\newtheorem{conjecture}{猜想}[section]

% 常用记号
\newcommand{\V}{\mathcal{V}}          % 不变量集合
\newcommand{\G}{\mathcal{G}}          % 图类
\newcommand{\Sn}{\mathcal{S}_n}       % n 点图样本
\newcommand{\F}{\mathcal{F}}          % 候选形式集合
\newcommand{\C}{\mathcal{C}}          % 候选猜想集合
\newcommand{\dd}{\mathrm{d}}
\newcommand{\R}{\mathbb{R}}
\newcommand{\N}{\mathbb{N}}
\newcommand{\ceil}[1]{\left\lceil #1 \right\rceil}
\newcommand{\floor}[1]{\left\lfloor #1 \right\rfloor}
\DeclareMathOperator*{\argmin}{arg\,min}
\DeclareMathOperator*{\argmax}{arg\,max}

% 章节标题样式
\titleformat{\section}
  {\normalfont\Large\bfseries\color{accent}}{\thesection}{0.8em}{}
\titleformat{\subsection}
  {\normalfont\large\bfseries\color{accent!85!black}}{\thesubsection}{0.7em}{}
\titlespacing*{\section}{0pt}{1.6ex plus .3ex}{1.0ex}
\titlespacing*{\subsection}{0pt}{1.2ex plus .2ex}{0.7ex}

% 页眉页脚
\pagestyle{fancy}
\setlength{\headheight}{24.5pt}
\fancyhf{}
\fancyhead[L]{\small\color{gray} 人工智能与数学 · 课程作业}
\fancyhead[R]{\small\color{gray} CRP-Agent}
\fancyfoot[C]{\small\thepage}
\renewcommand{\headrulewidth}{0.3pt}
\renewcommand{\headrule}{\hbox to\headwidth{\color{gray!45}\leaders\hrule height \headrulewidth\hfill}}

% 列表间距
\setlist[itemize]{leftmargin=1.6em, itemsep=2pt, topsep=3pt, parsep=0pt}
\setlist[enumerate]{leftmargin=1.8em, itemsep=2pt, topsep=3pt, parsep=0pt}

% 表格行距
\renewcommand{\arraystretch}{1.25}

% 图表标题样式：小号、居中、标签后用间隔而非冒号
\captionsetup{
  font=small,
  labelsep=quad,
  justification=centering,
  singlelinecheck=false
}

% 摘要环境（ctexart 自带 abstract，此处仅调整宽度）
\renewenvironment{abstract}{%
  \begin{center}\bfseries\large 摘\quad 要\end{center}
  \begin{quote}\small
}{%
  \end{quote}
}

% 强调框
\newtcolorbox{keybox}[1][]{
  enhanced, breakable,
  colback=accentlight, colframe=accent!60,
  boxrule=0.6pt, arc=2pt, left=8pt, right=8pt, top=6pt, bottom=6pt,
  fonttitle=\bfseries, title=#1
}

\newtcolorbox{notebox}[1][]{
  enhanced, breakable,
  colback=warnlight, colframe=orange!55!black,
  boxrule=0.6pt, arc=2pt, left=8pt, right=8pt, top=6pt, bottom=6pt,
  fonttitle=\bfseries, title=#1
}

% =====================================================================
\begin{document}

% ------------------------- 标题页 -------------------------
\begin{titlepage}
\centering
\vspace*{2.2cm}

{\color{accent}\rule{\linewidth}{1.2pt}}\\[0.5cm]
{\Huge\bfseries CRP-Agent}\\[0.35cm]
{\LARGE\bfseries 面向极值图论不等式的\\[0.2cm]
“猜想—反驳—证明”智能体}\\[0.5cm]
{\color{accent}\rule{\linewidth}{1.2pt}}

\vspace{1.6cm}

{\large\bfseries 课程作业报告}\\[0.3cm]
{\large 第一部分：研究问题}\\[1.8cm]

\begin{tabular}{rl}
\toprule
课程名称 & 人工智能与数学 \\
报告类型 & 阶段性报告（中期） \\
研究领域 & 组合数学 / 极值图论 · 数理统计 \\
撰写日期 & \today \\
\bottomrule
\end{tabular}

\vfill

{\small\color{gray}
本报告随研究推进持续修订。当前版本对应中期汇报阶段，\\[0.15cm]
仅包含问题陈述与技术路线，尚未包含实验数据与结论。}

\vspace{1.2cm}
\end{titlepage}

% ------------------------- 摘要 -------------------------
\pagenumbering{Roman}
\begin{abstract}
本报告陈述一项课程作业所研究的问题：\textbf{能否构建一个智能体，使其在极值图论的不变量不等式中自主地提出猜想、搜寻反例并尝试证明，同时以统计方法控制其产出的假阳性率？}

问题的出发点是一个观察：自 1980 年代 Fajtlowicz 的 Graffiti 以来，自动猜想系统已经能够批量生成图论不等式，但这类系统始终未能成为数学家的常规工具。障碍不在于“想不出猜想”，而在于\textbf{想出的猜想太多、太杂、太可能是偶然成立的}。当候选数量达到 $10^3$–$10^4$ 量级时，“哪些猜想值得投入人力去证明”本身成为一个多重检验问题，而现有系统无一系统性地处理它。

本文据此把研究问题拆成三个相互耦合的子问题：\emph{猜想生成}（在开放的形式空间中提出候选）、\emph{反例搜索}（以有限算力尽可能快地推翻候选）、\emph{统计筛选}（在大量检验中控制假发现率并分配搜索预算）。报告给出问题的形式化定义、与已有工作的区别、以及一条以“猜想—反驳—证明”（Conjecture–Refute–Prove, CRP）为骨架的智能体技术路线。

\vspace{0.6cm}
\noindent\textbf{关键词：}极值图论；图不变量；自动猜想；反例搜索；多重检验；假发现率；大语言模型智能体
\end{abstract}

\newpage
\tableofcontents
\newpage

\pagenumbering{arabic}
\setcounter{page}{1}

% ------------------------- 正文 -------------------------
% =====================================================================
%  第 1 节：研究问题的提出
% =====================================================================
\section{研究问题的提出}
\label{sec:problem}

\subsection{从一个具体的图论事实说起}
\label{sec:concrete}

设 $G=(V,E)$ 是一个有限简单无向图，$n=|V|$。记
\begin{itemize}
  \item $\alpha(G)$ 为\emph{独立数}：$G$ 中两两不相邻的顶点集的最大大小；
  \item $\omega(G)$ 为\emph{团数}：$G$ 中最大完全子图的大小；
  \item $\chi(G)$ 为\emph{色数}：使相邻顶点异色所需的最少颜色数。
\end{itemize}

这三个量满足一个初等的不等式：
\begin{equation}
  \chi(G) + \alpha(G) \;\le\; n + 1 .
  \label{eq:chi-alpha}
\end{equation}
证明只需一行：取最大独立集 $S$（$|S|=\alpha$），把 $S$ 整体染成一种颜色，其余 $n-\alpha$ 个顶点各染一种颜色，共需 $1+(n-\alpha)=n-\alpha+1$ 种颜色，故 $\chi\le n-\alpha+1$。

不等式 \eqref{eq:chi-alpha} 有两个值得注意的性质。第一，它是\textbf{紧的}：完全图 $K_n$ 上 $\chi=n,\ \alpha=1$，取等号；空图 $\overline{K_n}$ 上 $\chi=1,\ \alpha=n$，也取等号。第二，它的\textbf{证明只需要一行}，但它的\textbf{发现}并不显然——它把一个关于染色的问题（$\chi$）和一个关于独立集的问题（$\alpha$）通过顶点总数 $n$ 联系了起来。

图论中这类\textbf{不变量之间的不等式}数量极多，且其中相当一部分至今未被证明。例如

\begin{conjecture}[Reed, 1998]
\label{conj:reed}
对任意图 $G$，
\begin{equation}
  \chi(G) \;\le\; \ceil{\frac{\omega(G)+\Delta(G)+1}{2}},
  \label{eq:reed}
\end{equation}
其中 $\Delta(G)$ 为最大度。
\end{conjecture}

式 \eqref{eq:reed} 至今未被证明也未被否证。它的两个端点情形分别是已知定理：当 $\omega\le 2$ 时它退化为 Brooks 定理（$\chi\le\Delta+1$，除完全图与奇圈外），当 $\omega=\Delta+1$ 时它退化为平凡界 $\chi\le\Delta+1$。

\subsection{自动猜想：四十年的尝试与同一个瓶颈}
\label{sec:history-brief}

式 \eqref{eq:chi-alpha} 与式 \eqref{eq:reed} 形态相同：都是若干个图不变量的线性组合之间的不等式。一个自然的想法是：\textbf{既然这类不等式的形式如此规整，能否让机器自动地把它们找出来？}

这个想法已经被尝试了将近四十年。

\begin{itemize}
  \item \textbf{Graffiti}（Fajtlowicz，休斯敦大学，1986 年起）。第一个系统性的自动猜想程序。它从一个图数据库中提取不变量数值，用一套人工设计的启发式规则筛出“看起来成立”的不等式，累计提出逾千条猜想，其中若干条后来被数学家证明或否证。
  \item \textbf{Graffiti.pc}（DeLaViña，2000 年代）。Graffiti 的改进版本，调整了启发式以偏向“更可能可证”的猜想。
  \item \textbf{AutoGraphiX}（Caporossi \& Hansen，GERAD，2000 年代）。思路转向：不再从数据中“看”规律，而是用\textbf{优化}主动搜索使某个不等式最可能被违反的图，即主动构造反例；若搜索失败，则把该不等式作为猜想提出。
  \item \textbf{TxGraffiti}（Davila，2020 年代）。引入机器学习与线性规划，从数据中提取系数，并配套给出“极值图”作为猜想的证据。
\end{itemize}

四十年间，方法从启发式规则演进到优化、再到机器学习，但\textbf{有一个瓶颈始终没有变化}：

\begin{keybox}[核心观察]
这些系统的产出\textbf{远多于}人类能够处理的数量。Graffiti 提出逾千条猜想，其中被认真研究过的只是少数；其余的不是错的，而是\textbf{没人知道该不该花时间去证明}。

换句话说，自动猜想系统的困难不在于“生成猜想”——生成从来不是瓶颈——而在于\textbf{从海量候选中判断哪些值得投入}。
\end{keybox}

这个瓶颈有一个明确的技术根源。候选猜想是在\textbf{有限}的图样本上被检验的，而猜想本身断言的是\textbf{对所有图}成立。二者之间的鸿沟意味着：任何候选都可能是“在小图上碰巧成立”的假阳性。当候选数量达到 $10^3$–$10^4$ 量级时，即使每条候选的假阳性率很低，\textbf{期望的假阳性条数}仍然可能远超真实猜想条数。

于是问题变成了一个统计问题：\textbf{如何在大量并行检验中，控制“被接受的猜想里有多少是错的”？}这正是多重检验理论（multiple testing）的核心议题，而据我们所知，上述系统无一系统性地处理过它。

\subsection{本报告研究的问题}
\label{sec:question}

基于以上观察，本报告研究的\textbf{不是}某一个具体的图论猜想（如 Reed 猜想），而是一个方法层面的问题：

\begin{problem}[本课题的研究问题]
\label{prob:main}
能否构建一个智能体，使其在极值图论的不变量不等式中自主地完成“提出猜想—搜寻反例—尝试证明”的循环，并以统计方法控制其产出的假阳性率，使最终保留下来的猜想具备可辩护的可信度？
\end{problem}

这个问题包含三个相互耦合的子问题：

\begin{enumerate}[label=(\roman*), leftmargin=2.4em]
  \item \textbf{猜想生成}：给定一批图及其不变量数值，如何在\textbf{开放的形式空间}中提出候选不等式？（“开放”指：用哪些不变量、取线性还是非线性形式、是否含取整，这些都不预设。）
  \item \textbf{反例搜索}：给定一条候选，如何在\textbf{有限算力}下尽可能快地找到违反它的图？若找不到，如何说明“找不到”是有意义的而非搜索不充分？
  \item \textbf{统计筛选}：给定 $M$ 条候选及各自的检验结果，如何决定保留哪些？如何在候选之间\textbf{分配搜索预算}——把算力花在最有希望或最可疑的候选上？
\end{enumerate}

这三个子问题并非独立：生成策略决定了候选的分布，候选分布决定了筛选的难度，而筛选的反馈又应当指导下一轮生成。这种耦合正是需要一个\textbf{智能体}（而非一段脚本）来承载的原因——见第 \ref{sec:roadmap} 节。

\subsection{为什么选择这个问题}
\label{sec:why}

选择这一问题的理由有三条。

\paragraph{第一，它与统计学的训练直接相关。}
子问题 (iii) 是多重检验与最优实验设计的标准场景：$M$ 条候选对应 $M$ 次检验，需要在控制假发现率的前提下最大化真实发现数；预算分配则是 bandit 问题。这些工具在纯数学背景的研究者手中并不常用，但恰好是本课题的关键。这使课题具备\textbf{不可替代性}——它不是“用 AI 做数学”的又一个泛泛尝试，而是把统计方法论嵌入数学发现流程的一个具体切口。

\paragraph{第二，它满足课程对“体现智能体运行与迭代”的要求。}
第 \ref{sec:roadmap} 节将说明，CRP 循环的每一轮都会产生可记录、可比较、可审计的中间结果：提出的候选、执行的反例搜索、搜索的失败与成功、筛选的决策。这些记录构成完整的运行轨迹（trace），使得“智能体确实在迭代”这件事有据可查，而不必依赖最终是否解决了一个著名猜想。

\paragraph{第三，它的失败也是有信息量的。}
若智能体最终只“重新发现”了若干已知定理，这本身就是对系统正确性的验证；若它提出了在小图上成立、在更大图上被推翻的伪猜想，这些反例恰好量化了“小样本外推”的风险。相比之下，一个以“改进某个数值记录”为目标的课题，一旦数值上不去就几乎没有可汇报的内容。

\begin{notebox}[与课程要求的关系]
课程作业允许不彻底解决原始问题，但要求体现智能体的运行与迭代结果。本课题的“原始问题”（问题 \ref{prob:main}）是一个方法性问题，其价值不依赖于是否发现新的数学定理：即使最终只得到一个“能可靠地重现已知结论、并诚实报告其失败模式”的系统，也已经回答了一个有意义的问题——\emph{这类系统在什么条件下可信}。
\end{notebox}

% =====================================================================
%  第 2 章  已有工作
% =====================================================================

\section{已有工作}
\label{sec:background}

本节简要交代与本课题直接相关的工作，以说明本课题的位置，不展开技术细节。

\subsection{自动猜想程序}

最早系统性地从事该项工作的是 Fajtlowicz 于二十世纪八十年代开发的 Graffiti。该程序建立一个图数据库，为每张图计算各种不变量，再以一组人工设计的启发式规则挑选可能成立的不等式。据其作者统计，Graffiti 累计提出过上千条猜想，其中一部分后来被证明或推翻。

此后出现了若干改进版本，较知名的有 DeLaViña 的 Graffiti.pc、Caporossi 与 Hansen 的 AutoGraphiX，以及近年 Davila 的 TxGraffiti。这些工作在两个方向上有所加强：以线性规划代替手工规则确定系数，以及主动搜索反例而非被动检验。

上述系统的共同做法是产出候选、交由人工挑选。它们能够生成大量不等式，但均未将“如何从上千条候选中可靠地筛出少数几条”作为核心问题处理。

\subsection{基于大模型的方法}

2024 年以来，出现了将大语言模型接入数学发现流程的工作，代表性成果为 FunSearch 与 AlphaEvolve，以及 Charton 的 PatternBoost。这些工作在组合构造问题上取得了超过此前记录的结果。

这些工作处理的是构造问题而非猜想问题：其目标是找到一个具体的对象，而不是提出一条对全部对象成立的普遍命题。构造问题的优劣可由目标函数直接衡量，因而不需要处理候选过多、真假难辨的筛选问题。

本课题关注的正是后者。表\ref{tab:position} 给出两类工作与本课题的对比。

\begin{table}[H]
\centering
\small
\caption{相关工作与本课题的定位}
\label{tab:position}
\begin{tabular}{@{}lcc@{}}
\toprule
 & 是否处理候选过多 & 是否使用大语言模型 \\
\midrule
Graffiti 系列 & 否（人工挑选） & 否 \\
FunSearch 系列 & 不适用（构造型任务） & 是 \\
本课题 & 是（统计筛选） & 是 \\
\bottomrule
\end{tabular}
\end{table}

需要指出，上述空缺并非理论上的深刻发现，而更接近于一个被忽略的工程问题：研究多集中于提升生成的数量与质量，较少讨论生成之后如何处理。本课题的价值不在于提出新理论，而在于将这一问题认真处理一次，并给出可运行的实现与初步实证。

% =====================================================================
%  第 3 节：问题的形式化
% =====================================================================
\section{问题的形式化}
\label{sec:formulation}

本节把问题 \ref{prob:main} 写成可计算的形式。我们依次形式化\textbf{不变量向量}（\S\ref{sec:f-inv}）、\textbf{猜想空间}（\S\ref{sec:f-space}）、\textbf{反驳问题}（\S\ref{sec:f-refute}）、以及\textbf{筛选准则}（\S\ref{sec:f-screen}），最后给出整体目标（\S\ref{sec:f-obj}）。

\subsection{不变量向量与图样本}
\label{sec:f-inv}

固定一个有限的不变量集合 $\V=\{f_1,\dots,f_d\}$，其中每个 $f_i$ 是定义在图类 $\G$ 上的实值函数，且在同构下不变。对 $G\in\G$ 记
\begin{equation}
  \bm{f}(G) \;=\; \bigl(f_1(G),\dots,f_d(G)\bigr)\;\in\;\R^{d}.
  \label{eq:fvec}
\end{equation}

本课题取 $d=8$，$\V=\{\,n,\ m,\ \Delta,\ \alpha,\ \omega,\ \chi,\ \nu,\ \rho\,\}$（见表 \ref{tab:invariants}）。选择这八项的理由：前三项计算代价可忽略，后五项涵盖“组合—代数”两类视角，且它们之间的已知关系（如 $\chi\ge\omega$、$\omega\le\alpha$ 的补图对应）足够丰富，能为候选提供非平凡的约束。

\begin{definition}[图样本]
\label{def:sample}
称有限集合 $\Sn\subset\G_n$ 为一个\emph{$n$ 阶样本}，其中 $\G_n$ 为 $n$ 个顶点的所有简单无向图在同构下的等价类集合。记
\begin{equation}
  \bm{F}(\Sn) \;=\; \bigl\{\bm{f}(G) : G\in\Sn\bigr\}\;\subset\;\R^{d}
  \label{eq:sample}
\end{equation}
为样本的不变量矩阵（按行堆叠）。
\end{definition}

$|\G_n|$ 的增长极快：$|\G_7|=1044$，$|\G_8|=12{,}346$，$|\G_9|=274{,}668$，$|\G_{10}|=12{,}005{,}168$。因此 $n\le 9$ 可以\textbf{完全枚举}，$n\ge 10$ 必须\textbf{采样}。这一分界线是本课题实验设计的基本约束。

\subsection{猜想空间}
\label{sec:f-space}

\subsubsection{线性形式}

最自然的猜想形式是\textbf{线性不等式}：
\begin{equation}
  \bm{c}^{\top}\bm{f}(G) + c_0 \;\le\; 0 \qquad \text{对所有 } G\in\G,
  \label{eq:linform}
\end{equation}
其中 $\bm{c}\in\R^{d}$ 为不变量系数、$c_0\in\R$ 为常数项。

以式 \eqref{eq:chi-alpha} 为例。按 \eqref{eq:fvec} 的约定取 $\bm{f}=(n,m,\Delta,\alpha,\omega,\chi,\nu,\rho)$，则 $\chi+\alpha\le n+1$ 等价于
\begin{equation}
  -n + \alpha + \chi - 1 \;\le\; 0,
  \label{eq:chi-alpha-lin}
\end{equation}
即 $\bm{c}=(-1,\,0,\,0,\,1,\,0,\,1,\,0,\,0)$，$c_0=-1$。

\begin{remark}[为什么常数项必须显式出现]
\label{rem:const}
若想把 $c_0$ 并入某个不变量的系数，就需要把不等式除以该不变量，例如写成 $\bm{c}^\top\bm{f}/n + c_0/n\le 0$——这已不是线性形式。而 $n$ 在图之间是变化的，常数项无法被任何不变量的线性项吸收。

因此形式空间必须含 $c_0$，否则连式 \eqref{eq:chi-alpha} 这样的基本定理都无法表达。反过来，$c_0$ 也带来了麻烦：它是唯一\textbf{不能由样本直接定出}的自由度（LP 只会把它推到可行域的边界上），所以由 LP 得到的候选往往“恰好紧贴样本”，这本身就是过拟合的一种表现。
\end{remark}

为了让猜想\textbf{可读}，需要约束 $\bm{c}$ 的取值：系数应取小整数（如 $\{-3,\dots,3\}$），且非零项数不超过 $k$（如 $k\le 4$）。否则 LP 会给出系数为 $0.00314$ 的“猜想”，人类无法理解也无从证明。

\begin{definition}[候选形式]
\label{def:form}
一个\emph{候选形式}是一个三元组 $\varphi=(S,\sigma,k)$，其中
\begin{itemize}
  \item $S\subseteq\V$ 为参与的不变量集合，$|S|\le k$；
  \item $\sigma\in\{+,-\}$ 指示不等式方向（$\le$ 或 $\ge$，后者可通过乘以 $-1$ 归一化）；
  \item $k$ 为稀疏度上界。
\end{itemize}
记全体候选形式之集为 $\F$。
\end{definition}

$\F$ 的大小为 $\sum_{j\le k}\binom{d}{j}\cdot 2$，当 $d=8,k=4$ 时为 $2\cdot(8+28+56+70)=324$。这个规模是\textbf{可穷举}的，但穷举得到的候选缺乏“意图”——见 \S\ref{sec:f-llm}。

\subsubsection{系数求解：一个线性规划}

给定形式 $\varphi=(S,+,\cdot)$，问“是否存在系数使式 \eqref{eq:linform} 在样本 $\Sn$ 上成立”是一个线性规划可行性问题：
\begin{equation}
  \text{find } (\bm{c},c_0)\in\R^{|S|+1} \quad\text{s.t.}\quad
  \bm{c}^{\top}\bm{f}_{S}(G) + c_0 \le 0 \ \ \forall G\in\Sn,
  \label{eq:lp}
\end{equation}
其中 $\bm{f}_S$ 为 $\bm{f}$ 在 $S$ 上的投影。若可行域非空，还需从可行域中选一个“好”的 $(\bm{c},c_0)$——通常要求系数接近整数、范数小、非零项少。

注意 \eqref{eq:lp} 是\textbf{齐次化}的：把 $c_0$ 视为增广向量 $\tilde{\bm{c}}=(\bm{c},c_0)$ 与增广特征 $\tilde{\bm{f}}=(\bm{f}_S,1)$ 的内积，则约束变为 $\tilde{\bm{c}}^\top\tilde{\bm{f}}(G)\le0$。这是标准的齐次线性可行性问题，可用线性规划或（更强的）线性可分性判据求解。

\begin{remark}[平凡解的排除]
\label{rem:trivial}
式 \eqref{eq:lp} 总有平凡解 $\tilde{\bm{c}}=\bm{0}$。更麻烦的是\textbf{退化解}：若某不变量恒非负（如 $n\ge 1$、$\Delta\ge 0$），取 $\tilde{\bm{c}}$ 只在 $n$ 上为负、其余为 $0$ 即可得到“所有图都满足”的平凡猜想。

因此必须加入\textbf{非平凡性约束}，至少三条：
\begin{enumerate}[label=(\roman*)]
  \item 至少两个不变量系数非零，且\textbf{符号相反}（否则不等式只是“单调性”的重述）；
  \item 候选不能在全部样本上\textbf{取等号}（那说明它只是恒等式，不是不等式）；
  \item 候选的\textbf{紧性}不能过松：若 $\max_G \tilde{\bm{c}}^\top\tilde{\bm{f}}(G)$ 远小于 $0$，则该不等式虽真但无信息量（例如 $\chi\le 10n$ 恒成立却无用）。
\end{enumerate}
第 (iii) 条给出一个可操作的过滤阈值，但阈值的选取本身需要实验确定——过严会滤掉真正的不等式（它们往往不是最紧的），过松会放进大量无用候选。
\end{remark}

\begin{remark}[为什么不让 LLM 直接给系数]
\label{rem:llm-coef}
让 LLM 直接输出 $\bm{c}$ 是可行的，但有两个问题：其一，LLM 的数值直觉不可靠，给出的系数往往在样本上不成立，需要反复纠正；其二，即使成立，也难以判断是“真的紧”还是“系数取得过松”。把系数交给 LP，LLM 只负责形式，可使\textbf{可行性由求解器保证}，LLM 的错误表现为“形式不可行”这一干净信号，而非难以诊断的数值错误。
\end{remark}

\subsubsection{LLM 在猜想生成中的角色}
\label{sec:f-llm}

形式空间 $\F$ 虽小（324 个），但穷举有两个缺陷：

\begin{enumerate}[label=(\arabic*)]
  \item \textbf{大量形式必然不可行。}例如 $\{\alpha,\omega\}$ 上的 $\alpha\le c\,\omega$ 在任何含 $K_3$ 且 $\alpha$ 大的图上都失败。穷举会浪费算力在这些形式上。
  \item \textbf{穷举无法产生非线性形式。}形如 $\chi\le\ceil{(\omega+\Delta+1)/2}$ 的猜想含取整、含除法，不在 \eqref{eq:linform} 的范围内，但恰恰是这类形式最有数学意义。
\end{enumerate}

因此让 LLM 承担两件事：
\begin{itemize}
  \item \textbf{提议形式}：读入样本的统计摘要（各不变量的取值范围、极值图、已知定理清单），提议“值得一试”的形式，包括非线性形式；
  \item \textbf{解释失败}：当某形式不可行或候选被反例推翻时，读入反例的结构特征，说明“为什么失败”，并据此提议修正后的形式。
\end{itemize}

第二件事是循环得以\textbf{迭代}的关键：它不是重试，而是基于失败原因的定向修正。

\subsection{反驳问题}
\label{sec:f-refute}

给定候选 $(\bm{c},c_0)$，反驳即求解
\begin{equation}
  G^{\star} \;=\; \argmax_{G\in\G}\; \bigl(\bm{c}^{\top}\bm{f}(G)+c_0\bigr).
  \label{eq:refute}
\end{equation}
若 $\bm{c}^{\top}\bm{f}(G^{\star})+c_0>0$，则 $G^{\star}$ 是反例，候选被推翻；否则候选在 $\G$ 上（就搜索所及）成立。记
\begin{equation}
  \mathcal{E}(\bm{c},c_0) \;=\; \max_{G\in\G}\bigl(\bm{c}^{\top}\bm{f}(G)+c_0\bigr)
  \label{eq:excess}
\end{equation}
为候选的\textbf{超额量}：$\mathcal{E}>0$ 即存在反例，$\mathcal{E}\le0$ 即候选成立。

问题 \eqref{eq:refute} 是\textbf{无约束的离散优化}，$\G$ 无穷（$n$ 无上界）。实际求解需做三件事：

\paragraph{(a) 限定搜索域。}取 $n\le N_{\max}$（如 $N_{\max}=12$），并在 $n>9$ 时用采样代替枚举。这是必要的妥协，也是外推风险的来源。

\paragraph{(b) 设计搜索策略。}不同策略覆盖不同的“反例形态”：
\begin{itemize}
  \item \textbf{随机搜索}：$G(n,p)$ 上的随机图，$p$ 在若干值上扫描。适合发现“典型”反例。
  \item \textbf{结构化搜索}：正则图、二部图、Mycielski 构造（保持 $\omega$ 不变而增大 $\chi$）、Kneser 图、冠图（corona）。适合发现“极端”反例。
  \item \textbf{局部搜索}：从当前最优图出发做边翻转（edge flip），爬山或模拟退火。适合精化。
  \item \textbf{精确搜索}：$n\le 8$ 时直接枚举全部图。适合\textbf{确认}“小图上无反例”。
\end{itemize}

\paragraph{(c) 记录搜索的“充分性”。}若搜索未找到反例，需要说明搜索是否充分。我们记录每次搜索的\textbf{覆盖率}（枚举时的比例、采样时的样本量）与\textbf{最优违反值} $\max \bm{c}^\top\bm{f}$ 的分布。若最优值随预算增加而持续上升，说明搜索尚不充分；若趋于平台，则可初步认为无反例。

\begin{remark}[“无反例”不是“成立”]
\label{rem:no-counter}
问题 \eqref{eq:refute} 的求解只能给出“在搜索所及的范围内无反例”。这在逻辑上是弱结论，但在实践中是有信息量的：若在 $n\le 12$ 的数百万张图上、用四种策略搜索均未找到反例，则该候选至少\textbf{值得进一步研究}。本课题不声称能证明猜想，只声称能\textbf{排序}猜想。
\end{remark}

\subsection{统计筛选}
\label{sec:f-screen}

这是本课题的核心。设有 $M$ 条候选 $\mathcal{C}=\{c_1,\dots,c_M\}$，每条经若干轮反驳搜索后仍存活。问题是：\textbf{保留哪些？}

\subsubsection{困难：没有天然的 $p$ 值}

标准多重检验要求每条假设有一个在 $H_0$ 下服从已知分布的检验统计量。这里没有：
\begin{itemize}
  \item “未找到反例”是\textbf{确定性}结果，不含随机性；
  \item 但“未找到”依赖搜索预算，预算不同则结论强度不同；
  \item 且候选之间\textbf{相关}（共享不变量、共享样本），破坏独立性假设。
\end{itemize}

\subsubsection{方案：以“留出集上的违反裕度”为统计量}

我们构造如下统计量。把图样本分为三部分：
\begin{itemize}
  \item \textbf{发现集} $\mathcal{D}$：用于生成候选（$n\le 9$ 的枚举图）；
  \item \textbf{验证集} $\mathcal{V}$：用于初筛（$n=10,11$ 的采样图）；
  \item \textbf{测试集} $\mathcal{T}$：用于终检（$n=12$ 以上的采样图，含结构化图）。
\end{itemize}

对候选 $(\bm{c}_j,c_{0j})$ 定义其\textbf{违反裕度}
\begin{equation}
  T_j \;=\; \max_{G\in\mathcal{V}}\; \bigl(\bm{c}_j^{\top}\bm{f}(G)+c_{0j}\bigr),
  \label{eq:margin}
\end{equation}
即候选在验证集上的最大超额量。$T_j\le 0$ 表示在验证集上成立；$T_j$ 越接近 $0$，候选越“紧”，越可能在大图上被推翻。

\textbf{$T_j$ 的零分布可通过置换构造。}具体做法：保持 $\bm{c}_j$ 的\textbf{稀疏模式}（哪些不变量非零、各自符号）不变，在发现集上把不变量矩阵的各列随机置换（打破不变量之间的真实关联而保持各自的边缘分布），重新求解 LP \eqref{eq:lp} 得到 $\tilde{\bm{c}}_j,\tilde{c}_{0j}$，再在验证集上算出 $\tilde{T}_j$。重复 $B$ 次得到 $\{T_j^{(1)},\dots,T_j^{(B)}\}$，据此给出单侧 $p$ 值
\begin{equation}
  p_j \;=\; \frac{1+\#\{b: \tilde{T}_j^{(b)} \le T_j\}}{B+1}.
  \label{eq:pval}
\end{equation}

\begin{remark}[$p$ 值的方向]
\label{rem:pdir}
这里检验的原假设 $H_j$ 是“候选 $j$ 为假（存在反例）”。$p$ 值小意味着 $T_j$ 显著\textbf{低于}零分布——即候选在验证集上比“随机结构”更紧地成立，证据\textbf{反对} $H_j$，故拒绝 $H_j$、保留候选。若把不等号方向写反，$p$ 值会在候选\textbf{明显为假}时变小，BH 就会优先保留最可疑的候选——这是一个容易犯且不易察觉的错误。
\end{remark}

\begin{notebox}[这个构造是启发式的，需要验证]
置换检验的有效性依赖于“置换下的零分布确实对应 $H_0$”。在不变量数据上，这一点的严格性并不显然：不变量之间存在确定性关系（如 $\omega\le\chi\le\Delta+1$），随机置换列会破坏这些关系，产生的“伪图”不变量组合可能不对应任何真实图，使零分布偏离。

因此该方案\textbf{必须在实验中校准}：用已知定理（应当保留）检验其\textbf{功效}，用随机生成的无意义候选（不应当保留）检验其\textbf{水平}。若校准显示零分布明显偏离，备选方案是不做置换、改用“在不同 $n$ 的子样本间重采样”构造零分布。这是本课题的一个待验证环节。
\end{notebox}

\subsubsection{筛选规则}

有了 $p_j$，即可应用 BH 过程：
\begin{equation}
  k^{\star} = \max\left\{k : p_{(k)} \le \frac{k}{M}\,q\right\},
  \qquad \text{保留前 } k^{\star} \text{ 条}.
  \label{eq:bh}
\end{equation}
其中 $q$ 为目标 FDR 水平（如 $q=0.05$）。

但直接套用 BH 有两个问题，需要额外处理：
\begin{itemize}
  \item \textbf{相关性}：BH 在正相依下仍成立（Benjamini–Yekutieli 条件下），但不变量数据的相关性结构未知。我们计划用\textbf{BY 校正}（除以调和级数）作为保守版本，并报告两种结果。
  \item \textbf{预算不均}：不同候选的搜索预算不同，$p$ 值的可比性受影响。这引出下一小节。
\end{itemize}

\subsubsection{预算分配}

总预算 $B$ 需在 $M$ 条候选间分配。设候选 $j$ 在预算 $b$ 下被推翻的概率为 $\pi_j(b)$，未知且随 $b$ 递增。我们用\textbf{Beta–Bernoulli 汤普森采样}在线学习：

\begin{enumerate}[label=\arabic*.]
  \item 对每条候选维护后验 $\mathrm{Beta}(a_j,b_j)$，初值 $(1,1)$；
  \item 每轮按后验采样 $\theta_j\sim\mathrm{Beta}(a_j,b_j)$，把下一单位预算给 $\theta_j$ 最大的候选；
  \item 执行搜索，若推翻则记成功（$a_j\mathrel{+}=1$），否则记失败（$b_j\mathrel{+}=1$）；
  \item 重复至预算耗尽。
\end{enumerate}

这样做的效果是：可疑的候选获得更多预算，已被反复搜索仍无果的候选预算递减。\textbf{但这也引入偏差}——被少搜的候选，其“未找到反例”的证据更弱。因此预算分配与 $p$ 值计算必须联合设计：我们计划\textbf{按预算分层}计算 $p$ 值，或在最终筛选中对低预算候选施加惩罚。

\subsection{整体目标}
\label{sec:f-obj}

综合以上，本课题要优化的整体目标是：

\begin{keybox}[形式化目标]
给定不变量集合 $\V$、图样本 $\Sn$、总预算 $B$、目标 FDR 水平 $q$，求一个策略 $\pi$（由 LLM 提议形式、LP 定系数、搜索找反例、统计筛选四部分构成），使得在预算 $B$ 内：
\begin{equation}
  \max_{\pi}\; \mathbb{E}\bigl[\,|\mathcal{C}_{\text{retained}}|\,\bigr]
  \quad\text{s.t.}\quad
  \mathrm{FDR}\bigl(\mathcal{C}_{\text{retained}}\bigr)\le q,
  \label{eq:objective}
\end{equation}
其中 $\mathcal{C}_{\text{retained}}$ 为最终保留的猜想集合，FDR 按式 \eqref{eq:fdr} 定义，“假”指“存在反例”。
\end{keybox}

这个目标有三层含义，值得分开说明：

\begin{enumerate}[label=(\roman*)]
  \item \textbf{它是一个权衡，不是一个极大化。}提高保留数会推高 FDR，反之亦然。策略的优劣体现在\textbf{帕累托前沿}上：在给定 $q$ 下能保留多少条。
  \item \textbf{它可被部分评估。}虽然“是否真的成立”无法完全判定，但可以用\textbf{已知定理}（真阳性）与\textbf{已知反例}（真阴性）构造一个有标签的子集，在该子集上评估策略的召回与精确率。这是可行的验证路径。
  \item \textbf{它的价值在于过程。}即使最终保留数为零（所有候选都被推翻或无法通过筛选），循环中产生的反例、失败模式、以及“为什么筛不出来”的诊断，本身就是对问题 \ref{prob:main} 的回答。
\end{enumerate}

% =====================================================================
%  第 4 节：问题为何未解决
% =====================================================================
\section{问题为何未解决}
\label{sec:openness}

问题 \ref{prob:main} 不是一个已被解决的问题，也不是一个“工程实现问题”。本节说明它难在哪里——分为\textbf{数学层面的困难}（\S\ref{sec:hard-math}）、\textbf{统计层面的困难}（\S\ref{sec:hard-stat}）、以及\textbf{方法论层面的困难}（\S\ref{sec:hard-method}）。把困难讲清楚，也是说明本课题为何“允许不彻底解决”的前提。

\subsection{数学层面的困难}
\label{sec:hard-math}

\subsubsection{候选空间的开放性与可读性的矛盾}

若不加限制，猜想空间是无穷的：可以引入任意复合不变量（如 $\chi^2/\omega$、$\log\alpha$、$\Delta^{\alpha}$），可以加入任意取整与条件分支。理论上，越大的空间越可能包含“有趣的”猜想，但同时也使候选的可读性与可证性急剧下降。

这里存在一个\textbf{实质性的张力}：

\begin{itemize}
  \item 限制在\textbf{线性、小整数系数、稀疏}的形式内，候选可读、可用 LP 求解、可用初等方法验证——但已知的大量有趣猜想（Reed 猜想的取整、Nordhaus–Gaddum 型的两图关系）落在这个空间之外。
  \item 放开到\textbf{非线性、含取整}的形式，表达力足够——但“求解系数”不再是 LP，而变成混合整数非线性规划，且“这个猜想是否已知”的判断也变得困难。
\end{itemize}

本课题采取折中：\textbf{以线性形式为主干（可自动化），用 LLM 提议的非线性形式作为补充（半自动化）}。非线性候选不进入 LP，而由 LLM 给出系数后直接检验。这使系统能触及 Reed 猜想这类形式，但也意味着这部分候选的质量控制更依赖 LLM，需要单独评估。

\subsubsection{“已知”的判定问题}

一条候选是否“新”，没有算法化的判据。数学文献中的等价变形（移项、乘以正数、代入已知不等式）可以产生无穷多种表述同一事实的方式。判断“$\chi+\alpha\le n+1$”与“$\alpha\le n-\chi+1$”是同一个定理，需要语义层面的理解。

这构成一个\textbf{无法完全自动化}的环节。本课题的做法是：维护一个已知结论的清单（数十条经典不等式），用 LLM 做语义匹配，并\textbf{明确标注}“新颖性判断是启发式的、可能出错”。在最终报告中，所有“候选新猜想”都会附上人工核对的结果。

\subsubsection{反例的存在性不可判定}

对一般的 $\bm{c}$，判定“是否存在图使 $\bm{c}^\top\bm{f}>0$”在理论上是困难的：$\G$ 无穷，且 $\alpha,\omega,\chi$ 的计算本身是 NP-难的。即使限定 $n\le N$，枚举 $\G_N$ 的代价也随 $N$ 指数增长（$|\G_{10}|\approx1.2\times10^7$，$|\G_{11}|\approx10^9$）。

因此在实践中，反驳永远是\textbf{不完整的}：我们只能在一个可及的范围内搜索。这与“证明”有本质区别——证明是对全体的断言，搜索只是对部分的检查。

\subsection{统计层面的困难}
\label{sec:hard-stat}

\subsubsection{零分布未知}

\S\ref{sec:f-screen} 提出的置换检验方案，其有效性依赖于一个未经验证的假设：置换后得到的 $\tilde{T}_j$ 分布确实近似于 $H_0$ 下的分布。问题在于：

\begin{itemize}
  \item 不变量之间存在\textbf{确定性约束}（$\omega\le\chi$、$\chi\le\Delta+1$、$\alpha\omega\ge n$ 等）。随机置换会破坏这些约束，产生的“伪图”不变量组合可能不对应任何真实图，从而使零分布偏离真实情形。
  \item 候选之间\textbf{强相关}。同一个不变量出现在多条候选中，一条被推翻往往意味着相关的几条也被推翻。BH 在任意相依下不保证 FDR 控制。
\end{itemize}

这些问题的严重程度\textbf{只能通过实验评估}。我们计划用“已知定理 + 已知反例”构成的有标签集合校准：若某策略在该集合上的 FDR 显著超过名义水平 $q$，则说明零分布构造有误，需要修正（例如改用 \S\ref{sec:f-refute} 中记录的“搜索充分性”作为协变量做条件检验）。

\subsubsection{预算分配引入的选择偏差}

汤普森采样把预算集中到“看起来可疑”的候选上，这本身是合理的——但它使不同候选的“未被推翻”具有\textbf{不同的证据强度}。被少搜的候选，其 $T_j$ 更可能是“没搜到”而非“真的没有”。

这是一个典型的\textbf{选择性推断}（selective inference）问题：数据依赖的决策（把预算给谁）破坏了后续检验的分布性质。文献中的处理方式包括样本分割（用一部分预算决定分配、另一部分做检验）与条件推断（在给定分配的条件下做检验）。本课题倾向于\textbf{样本分割}，因为它简单且不需要推导复杂的条件分布，代价是浪费一部分预算。

\subsubsection{外推的系统性偏移}

$n\le 9$ 的图与 $n=12$ 的图不只是“大小不同”，而是\textbf{结构上不同}。例如：

\begin{itemize}
  \item 需要较多顶点才能实现的构造（如某些 Mycielski 迭代、Kneser 图 $K(5,2)$ 需 10 个顶点）在 $n\le 9$ 中完全缺席；
  \item 小图上“几乎所有图都是连通的、度分布集中”，而大图上存在更多极端结构；
  \item 某些不等式在小图上因“没有足够的空间构造反例”而成立，在大图上则失效。
\end{itemize}

因此从发现集到测试集的外推，其偏移是\textbf{系统性}的，不能靠增大样本量消除。这构成本课题最根本的方法论限制，也是我们在结论中必须明确声明的。

\subsection{方法论层面的困难}
\label{sec:hard-method}

\subsubsection{“成功”的定义本身是模糊的}

一个自动猜想系统怎样算成功？可能的判据包括：

\begin{itemize}
  \item 提出的猜想中有多少后来被证明（需要数年时间验证，无法在课程周期内评估）；
  \item 提出的猜想中有多少“有意思”（主观）；
  \item 系统能否重现已知定理（可评估，但重现已知不算发现）；
  \item 系统能否产出可被人类数学家在一小时内判断真伪的猜想（可评估，且实用）。
\end{itemize}

本课题采用最后一条作为主要判据，辅以第三条作为正确性检验。这个选择本身是一个\textbf{方法论主张}：自动猜想系统的价值在于\textbf{降低人类筛选成本}，而不在于替代人类做证明。

\subsubsection{LLM 的不稳定性}

LLM 的输出具有随机性：同一个提示在不同时刻可能给出不同的形式提议。这带来两个问题：

\begin{itemize}
  \item \textbf{可复现性}：整个循环的结果依赖于 LLM 的采样。我们通过记录每次调用的完整输入输出（写入轨迹文件）、并固定随机种子与模型版本来缓解，但无法完全消除。
  \item \textbf{评估的混淆}：若最终结果不理想，难以判断是“方法不行”还是“这次 LLM 恰好表现差”。需要通过多次独立运行（不同种子）报告结果的分布，而非单次结果。
\end{itemize}

\subsubsection{“智能体”与“脚本”的界限}

这是本课题必须自我审视的一点。一个自然的质疑是：既然猜想生成可以用 LP、反驳可以用搜索、筛选可以用 BH，那么 LLM 在其中是否只是装饰？若去掉 LLM 系统仍能运行，是否说明它不构成“智能体”？

我们的回应是：LLM 承担的是\textbf{在开放空间中的定向搜索}——提议哪些形式值得试、根据失败原因修正形式、判断候选是否已知。这些环节的规则无法预先枚举（否则就退化为穷举），需要语义理解。但这一主张需要\textbf{实验支持}：我们计划做消融实验，比较“有 LLM 提议”与“随机/穷举提议”在相同预算下的产出质量。若两者无显著差异，则应诚实报告“LLM 在此任务中贡献有限”——这同样是一个有价值的结论。

\begin{notebox}[关于本课题的定位]
上述困难说明：本课题\textbf{不可能}在课程周期内被彻底解决。它的合理目标是——搭建一个可运行的 CRP 循环，在一个受限但明确的设定下（线性形式、$n\le 12$、已知不变量集合）产出可审计的结果，并对“统计筛选是否有效”给出初步的实证回答。这符合课程“允许不彻底解决、但须体现运行与迭代”的要求。
\end{notebox}

% =====================================================================
%  第 5 节：技术路线
% =====================================================================
\section{技术路线}
\label{sec:roadmap}

本节给出 CRP 循环的整体设计（\S\ref{sec:r-arch}）、四个模块的具体职责（\S\ref{sec:r-modules}）、以及可审计轨迹的设计（\S\ref{sec:r-trace}）。本节描述\textbf{计划}，实验结果将在后续版本中补充。

\subsection{整体架构}
\label{sec:r-arch}

\begin{figure}[htbp]
\centering
\begin{tikzpicture}[
  x=1cm, y=1cm,
  box/.style={draw=accent!70, fill=accentlight, rounded corners=2pt,
              minimum width=3.2cm, minimum height=0.95cm, align=center,
              font=\small, inner sep=3pt},
  llmbox/.style={box, fill=orange!12, draw=orange!60!black},
  statbox/.style={box, fill=green!10, draw=green!45!black},
  arr/.style={-{Stealth[length=2.0mm]}, thick, accent!85},
  fb/.style={-{Stealth[length=2.0mm]}, thick, dashed, orange!75!black},
  lbl/.style={font=\scriptsize\itshape, gray}
]

% ---- 第 1 列：猜想 ----
\node[llmbox] (llm)  at (0,0)     {LLM 提议形式\\[-1pt]{\scriptsize 读样本摘要 + 已知定理}};
\node[box]    (lp)   at (0,-1.55) {线性规划定系数\\[-1pt]{\scriptsize 可行性 + 整数化}};
\node[box]    (cand) at (0,-3.10) {候选猜想集 $\mathcal{C}$\\[-1pt]{\scriptsize 形式 + 系数 + 紧性}};

% ---- 第 2 列：反驳 ----
\node[box] (search) at (4.9,0)     {反例搜索\\[-1pt]{\scriptsize 随机/结构/局部/精确}};
\node[box] (budget) at (4.9,-1.55) {预算分配\\[-1pt]{\scriptsize Beta–Bernoulli 汤普森}};

% ---- 第 3 列：筛选 ----
\node[statbox] (margin) at (9.8,0)     {违反裕度 $T_j$\\[-1pt]{\scriptsize 验证集上最大违反量}};
\node[statbox] (perm)   at (9.8,-1.55) {置换零分布\\[-1pt]{\scriptsize $\to$ $p$ 值}};
\node[statbox] (bh)     at (9.8,-3.10) {BH / BY 筛选\\[-1pt]{\scriptsize FDR $\le q$}};

% ---- 阶段标注 ----
\node[font=\small\bfseries, accent] at (0,0.95)     {(1)\ 猜想};
\node[font=\small\bfseries, accent] at (4.9,0.95)   {(2)\ 反驳};
\node[font=\small\bfseries, accent] at (9.8,0.95)   {(3)\ 筛选};

% ---- 数据流 ----
\draw[arr] (llm) -- (lp);
\draw[arr] (lp) -- (cand);
\draw[arr] (cand.east) -- ++(0.55,0) |- (search.west);
\draw[arr] (budget.north) -- (search.south);
\draw[arr] (search) -- (margin);
\draw[arr] (margin) -- (perm);
\draw[arr] (perm) -- (bh);

% ---- 反馈回路（绕行图外，避免穿越方框）----
\draw[fb] (bh.south) -- (9.8,-4.45) -- (-2.45,-4.45) -- (-2.45,0) -- (llm.west);
\node[lbl, anchor=north] at (4.0,-4.55) {失败原因 + 反例结构 $\Rightarrow$ 修正下一轮的形式提议};

\end{tikzpicture}
\caption{CRP 循环的架构。实线为数据流，虚线为反馈。左侧“猜想”与右侧“筛选”通过反馈回路耦合：筛选阶段暴露的失败模式被送回 LLM，用于修正下一轮的形式提议。}
\label{fig:arch}
\end{figure}

图 \ref{fig:arch} 的三列对应问题 \ref{prob:main} 的三个子问题。值得注意的是\textbf{反馈回路}：筛选不只是终点，它产生的“为什么这条候选不可信”的信息会回馈到生成端。这是系统被称为“智能体”而非“流水线”的实质原因——它有一个闭合的、基于结果调整策略的循环。

\subsection{模块职责}
\label{sec:r-modules}

\subsubsection{模块一：图样本与不变量计算}

\textbf{输入}：$n$ 的上界 $N_{\max}$、采样策略。
\textbf{输出}：不变量矩阵 $\bm{F}\in\R^{|\Sn|\times d}$ 及对应的图（以边表存储）。

关键决策是\textbf{如何处理 NP-难的不变量}。我们采取三级策略：

\begin{enumerate}[label=(\arabic*)]
  \item \textbf{$n\le 8$ 精确计算。}$|\G_8|=12{,}346$，用带剪枝的精确算法（$\alpha$ 用分支定界，$\chi$ 用 DSATUR + 回溯）可在分钟级完成。
  \item \textbf{$n=9$ 用界代替。}$|\G_9|=274{,}668$，精确计算色数不可行。改用区间 $[\omega,\ \Delta+1]$，并记录区间宽度作为“不确定度”。$\alpha$ 用精确算法（9 点图的独立数计算很快）。
  \item \textbf{$n\ge 10$ 采样。}用 $G(n,p)$（$p$ 扫描）、正则图、结构化图三类采样，每类若干千张。
\end{enumerate}

\begin{notebox}[一个需要报告的取舍]
用区间代替精确色数，会使“$\chi$ 参与的候选”的判定变得保守：若 $\chi$ 的真值在区间内变动能改变候选的真假，则该候选被标记为“不确定”而非“成立”或“被推翻”。这会降低召回率，但避免了假阳性。这个取舍将在实验部分量化。
\end{notebox}

\subsubsection{模块二：猜想生成}

\textbf{输入}：不变量矩阵、已知定理清单。
\textbf{输出}：候选集合 $\mathcal{C}$，每条含形式、系数、在发现集上的紧性。

流程：
\begin{enumerate}[label=\arabic*.]
  \item 计算样本摘要：各不变量的取值范围、分位数、极值图、两两相关系数；
  \item 把摘要与已知定理清单送入 LLM，要求其输出若干候选形式（JSON 格式，含不变量子集、方向、是否含取整）；
  \item 对每个线性形式，求解 LP \eqref{eq:lp}；不可行则记录并跳过；
  \item 对可行的解，做整数化（把系数缩放到小整数，检验是否仍可行）；
  \item 计算每条候选在发现集上的\textbf{紧性} $\max_G \bm{c}^\top\bm{f}(G)$，紧性接近 $0$ 的候选优先进入反驳。
\end{enumerate}

\subsubsection{模块三：反例搜索}

\textbf{输入}：候选 $\bm{c}$、预算 $b$。
\textbf{输出}：反例（若找到）或“未找到”及搜索充分性指标。

四种策略按预算比例混合执行，比例由上一轮的成败调整：

\begin{table}[H]
\centering
\caption{反例搜索的四种策略及其覆盖的反例形态。}
\label{tab:search}
\small
\begin{tabularx}{\linewidth}{@{}p{1.9cm} >{\raggedright\arraybackslash}X >{\raggedright\arraybackslash}p{4.3cm}@{}}
\toprule
策略 & 做法 & 擅长发现的反例 \\
\midrule
随机 & $G(n,p)$，$p\in\{0.1,\dots,0.9\}$，每个 $n$ 若干样本 & “典型”反例，$p$ 接近临界值时 \\
结构化 & 正则图、二部图、Mycielski 构造、Kneser 图、冠图、完全多部图 & 极端反例（$\chi$ 大而 $\omega$ 小等） \\
局部 & 从当前最优图出发做边翻转，爬山或模拟退火 & 精化：把违反量从接近 $0$ 推到正 \\
精确 & $n\le 8$ 时全枚举 & \textbf{确认}小图上无反例 \\
\bottomrule
\end{tabularx}
\end{table}

\subsubsection{模块四：统计筛选}

\textbf{输入}：存活候选及其在各集合上的违反裕度。
\textbf{输出}：保留集合、每条候选的 $p$ 值与 FDR 估计。

按 \S\ref{sec:f-screen} 的方案执行，另加两项：
\begin{itemize}
  \item \textbf{校准}：用已知定理（应被保留）与已知反例（应被推翻）构成的有标签子集，检验筛选的召回率与精确率；
  \item \textbf{敏感性分析}：对 $q$、置换次数 $B$、预算分配策略做网格扫描，报告结果对超参数的依赖。
\end{itemize}

\subsection{可审计轨迹}
\label{sec:r-trace}

整个循环的每一步写入 JSONL 文件，每行一条记录，字段包括：

\begin{table}[H]
\centering
\caption{轨迹记录的主要字段。}
\label{tab:trace}
\small
\begin{tabular}{@{}llp{7.2cm}@{}}
\toprule
字段 & 类型 & 含义 \\
\midrule
\texttt{round} & int & 循环轮次 \\
\texttt{stage} & str & \texttt{propose} / \texttt{refute} / \texttt{screen} \\
\texttt{candidate\_id} & str & 候选标识（形式的哈希） \\
\texttt{form} & object & 形式：不变量子集、方向、是否含取整 \\
\texttt{coeffs} & list & LP 求得的系数 \\
\texttt{llm\_input} & str & 送入 LLM 的完整提示 \\
\texttt{llm\_output} & str & LLM 的原始输出 \\
\texttt{search\_strategy} & str & 本轮使用的搜索策略 \\
\texttt{search\_budget} & int & 消耗的预算（评估次数） \\
\texttt{counterexample} & object & 反例的边表与不变量值（若有） \\
\texttt{margin} & float & 违反裕度 $T_j$ \\
\texttt{p\_value} & float & 置换检验得到的 $p$ 值 \\
\texttt{decision} & str & \texttt{keep} / \texttt{drop} / \texttt{uncertain} \\
\bottomrule
\end{tabular}
\end{table}

轨迹的作用有三：
\begin{enumerate}[label=(\arabic*)]
  \item \textbf{可复现}：给定轨迹，可重放整个循环（LLM 输出已记录，不需重新调用）；
  \item \textbf{可诊断}：分析“哪些形式的候选最容易被推翻”“哪类搜索策略最有效”；
  \item \textbf{可汇报}：中期与最终报告的图表直接由轨迹生成，避免手工整理。
\end{enumerate}

\subsection{验证计划}
\label{sec:r-valid}

由于“是否真的成立”无法完全判定，验证分三层：

\begin{enumerate}[label=(\roman*)]
  \item \textbf{正确性检验}：系统能否重现一批已知定理（Brooks、$\chi+\alpha\le n+1$、$\alpha\ge n/(\Delta+1)$ 等）？能否对已知的反例图（如 Grötzsch 图对某些断言）正确判定？
  \item \textbf{筛选有效性}：在有标签子集上，BH 筛选的实际 FDR 是否接近名义水平 $q$？与不做筛选（全保留）相比，精确率提升多少？
  \item \textbf{消融}：LLM 提议 vs.\ 随机提议 vs.\ 穷举提议，在相同预算下的产出质量差异；有预算分配 vs.\ 均匀分配。
\end{enumerate}

第 (i) 层是“系统能不能用”，第 (ii) 层是“统计方法有没有用”，第 (iii) 层是“智能体成分有没有用”。三层都通过，才能支持“本课题提出了一个有效方法”的主张；任一层失败，也是有价值的负面结果。

\subsection{当前进度与后续安排}
\label{sec:r-progress}

\begin{table}[H]
\centering
\caption{工作安排。表中“状态”截至本报告撰写时。}
\label{tab:schedule}
\small
\begin{tabular}{@{}clp{6.4cm}l@{}}
\toprule
阶段 & 内容 & 产出 & 状态 \\
\midrule
D0 & 选题论证、报告框架 & 本报告第一部分 & 完成 \\
D1 & 图枚举与不变量计算 & 不变量矩阵（$n\le 9$） & 待办 \\
D2 & LP 猜想生成器 & 首批候选 & 待办 \\
D3 & 反例搜索器 & 首轮反驳结果 & 待办 \\
D4 & 统计筛选与校准 & FDR 分析 & 待办 \\
D5 & 案例研究与消融 & 三条候选的深度分析 & 待办 \\
D6 & 中期汇报材料 & 幻灯片 & 待办 \\
\bottomrule
\end{tabular}
\end{table}


% ------------------------- 参考文献 -------------------------
\newpage
% =====================================================================
%  参考文献
% =====================================================================
\section*{参考文献}
\addcontentsline{toc}{section}{参考文献}
\markboth{参考文献}{参考文献}

\begin{sloppypar}
\begin{enumerate}[label={[\arabic*]}, leftmargin=2.6em, itemsep=3pt]

\item Fajtlowicz, S. On conjectures of Graffiti. \emph{Discrete Mathematics}, 1988, 72(1--3): 113--118.

\item Fajtlowicz, S. On conjectures of Graffiti, II. \emph{Congressus Numerantium}, 1988, 60: 187--197.

\item DeLaViña, E. Graffiti.pc: a variant of Graffiti. \emph{DIMACS Series in Discrete Mathematics and Theoretical Computer Science}, 2002, 69: 71--79.

\item Caporossi, G., Hansen, P. Variable neighborhood search for extremal graphs: 1. The AutoGraphiX system. \emph{Discrete Mathematics}, 2000, 212(1--2): 29--44.

\item Davila, R. R. TxGraffiti: a system for automated conjecturing in graph theory. 预印本, 2020 年代.

\item Reed, B. $\omega$, $\Delta$ and $\chi$. \emph{Journal of Graph Theory}, 1998, 27(4): 177--212.

\item Benjamini, Y., Hochberg, Y. Controlling the false discovery rate: a practical and powerful approach to multiple testing. \emph{Journal of the Royal Statistical Society: Series B}, 1995, 57(1): 289--300.

\item Benjamini, Y., Yekutieli, D. The control of the false discovery rate in multiple testing under dependency. \emph{The Annals of Statistics}, 2001, 29(4): 1165--1188.

\item Romera-Paredes, B., Barekatain, M., Novikov, A., et al. Mathematical discoveries from program search with large language models. \emph{Nature}, 2024, 625: 468--475.

\item Charton, F., Ellenberg, J., Wagner, A. Z., Williamson, G. PatternBoost: constrained optimization with transformers for constructive problems. \emph{arXiv:2411.19826}, 2024.

\item Novikov, A., Vũ, M. N., et al. AlphaEvolve: a Gemini-powered coding agent for designing advanced algorithms. DeepMind 技术报告, 2025.

\item Diestel, R. \emph{Graph Theory} (5th ed.). Springer, 2017.

\item Bollobás, B. \emph{Extremal Graph Theory}. Academic Press, 1978.

\item West, D. B. \emph{Introduction to Graph Theory} (2nd ed.). Prentice Hall, 2001.

\item Harary, F. \emph{Graph Theory}. Addison-Wesley, 1969.

\item Efron, B., Tibshirani, R. J. \emph{An Introduction to the Bootstrap}. Chapman \& Hall, 1993.

\item Good, P. \emph{Permutation, Parametric, and Bootstrap Tests of Hypotheses} (3rd ed.). Springer, 2005.

\item Thompson, W. R. On the likelihood that one unknown probability exceeds another in view of the evidence of two samples. \emph{Biometrika}, 1933, 25(3--4): 285--294.

\item Russo, D. J., Van Roy, B., Kazerouni, A., Osband, I., Wen, Z. A tutorial on Thompson sampling. \emph{Foundations and Trends in Machine Learning}, 2018, 11(1): 1--96.

\item Taylor, J., Tibshirani, R. J. \emph{Statistical Learning with Sparsity}. Chapman \& Hall, 2015.

\item Wilf, H. S. The eigenvalues of a graph and its chromatic number. \emph{Journal of the London Mathematical Society}, 1967, 42: 330--332.

\end{enumerate}
\end{sloppypar}


\end{document}
